P01 & What is conditioned on, and which identity unit should be preserved?\par\emph{Task/reference definition developed} & Donor experiments distinguish person, eye, side and acquisition. Inpainting fixes visible context and permits observed-reference tests where a real region is hidden.\par\textbf{Missing:} A model-specific conditioning contract and an explicit treatment of unidentifiable detail. Human anatomical relationships need domain-specific replacements outside the current data. \\
P02 & Which target strata, donor cohorts and acquisition domains are being compared?\par\emph{Donor matching available; generator strata not measured} & Age/recorded-sex-matched means and family exclusion define reference cohorts. Accessible face tier-1 outputs do not establish a population-stratified audit of generated completions.\par\textbf{Missing:} Prespecified target groups, metadata availability, independent-unit counts and matched task-difficulty strata. A modality or dataset contrast is not a demographic disparity. \\
P03 & Is the distance implementation and support contract explicit?\par\emph{Substantial saved implementation-probe coverage} & Validity-mask tests and saved official inpainting-evaluator probes examine tensor layout, clipping, normalization and within-case ordering.\par\textbf{Missing:} Deployment parity and task-specific mask-boundary/three-dimensional conventions. A metric name alone does not specify an evaluator. \\
P04 & Does the score notice attenuation toward a global or matched population mean?\par\emph{Mechanism evidence; MRI provisional} & Known-target synthetic shrinkage and saved donor profiles separate low reference error from retained residual amplitude and target direction.\par\textbf{Missing:} Actual generated completions, information recoverable from their context and a distinction between legitimate mean prediction and harmful loss of relevant variation. \\
P05 & Can omitted detail be distinguished from plausible but incorrect detail?\par\emph{Controlled proxy evidence} & Norm-matched incorrect residuals, means and noise stress-test amplitude and alignment separately. Whole-image ArcFace evaluates face donor identity, not an inpainted composite.\par\textbf{Missing:} Paired model/control comparisons at the same spatial scope, independent anatomical error assessments and a defined cost of omission versus incorrect structure. \\
P06 & Is acquisition variation separated from identity or anatomy?\par\emph{Partial; MRI gate unresolved} & Repeat captures, paired contrast sensitivity and synthetic pupil changes supply nuisance controls. HACA3 contrast reduction does not demonstrate preservation of individual anatomy.\par\textbf{Missing:} Accepted correspondence/QC where needed and task-specific checks that nuisance invariance retains meaningful anatomy. \\
P07 & Are anatomical structure, downstream utility and spatial support tested?\par\emph{Synthetic/partial spatial evidence} & Threshold-segmentation and topology toys, Gaussian maps and historical size/shell sweeps provide mechanism and spatial diagnostics.\par\textbf{Missing:} Independent real-image structure/usefulness criteria, irregular masks and completed object-relative scale studies. Downstream agreement alone is not identity validation. \\
P08 & How does a generator compare with the ladder within each target stratum?\par\emph{Proposed readout; no subgroup generator experiment} & Existing controls provide candidate reference levels for paired comparisons. No valid shared-target, shared-mask model-output group profiles are present in the accessible snapshot.\par\textbf{Missing:} Frozen model outputs, matched eligible targets/support, sampling budgets and paired output-minus-control profiles for each metric, rung and target group. \\
P09 & Do apparent group differences survive a common reference and comparable task difficulty?\par\emph{Proposed comparison; not measured} & Global and metadata-matched donor means exist as distinct comparators. They motivate retaining both raw distances and common-reference versus within-group coordinates.\par\textbf{Missing:} Comparable mask/location/acquisition conditions, uncertainty in reference pools and group differences, and sensitivity to group-specific reference variance. Descriptive gaps are not causal fairness findings. \\
P10 & Are between-target variation and within-context sampling variation separated?\par\emph{Conceptual controls only} & Averaging and attenuation motivate the distinction, but saved donor averages are not samples from a trained generator.\par\textbf{Missing:} Repeated outputs at a fixed seed/sample budget, conditional diversity and target-specific fidelity measurements. Reduced variance or mean preference alone does not establish mode collapse. \\
P11 & Do profiles support stable, independently meaningful evaluation decisions?\par\emph{Metric dependence partial; model/group decisions unvalidated} & Rung correlations, source intervals and toy rank probes exist. They are not confidence intervals for method rankings or group disparities.\par\textbf{Missing:} Independent quality criteria, target-cluster uncertainty with repeated masks/seeds, eligibility/missingness and worst-group summaries where estimable, and frozen aggregation choices. \\
P12 & Which parts transfer beyond inpainting to other generative settings?\par\emph{Proposed extension; human-anatomy evidence only} & Inpainting exposes context, withheld content and population controls in a paired task. These components can define an audit where a broader generator has an identity-bearing condition and meaningful reference unit.\par\textbf{Missing:} New domain-specific controls and validation. Unconditional generation has no intended target individual; paired identity fidelity cannot be applied without that condition and requires separate distribution/coverage criteria. \\
\bottomrule
